% =============================================================================
% Chapter 00 — Airbnb CDMX Dashboard
% =============================================================================

\chapter{Airbnb CDMX Dashboard}
\label{ch:airbnb}

\sourcefiles{%
  \filepath{projects/00-demo-aestehtics/data-pipeline/airbnb\_etl.py},
  \filepath{projects/00-demo-aestehtics/src/app/airbnb/page.tsx},
  \filepath{projects/00-demo-aestehtics/src/components/AirbnbDashboard.tsx},
  \filepath{projects/00-demo-aestehtics/src/components/charts/},
  \filepath{projects/00-demo-aestehtics/src/app/globals.css}
}

This chapter covers the lightest project in the portfolio in terms of statistical
complexity, but one of the most instructive in terms of frontend architecture and
visualization design decisions. The Airbnb CDMX dashboard is a Next.js~14
application that presents Mexico City's short-term rental market to a non-technical
audience using five chart types, a dark/light mode system, and a static-data
architecture that requires no backend at runtime.

% ---------------------------------------------------------------------------
\section{Business Context}
\label{sec:airbnb-business}
% ---------------------------------------------------------------------------

Mexico City (CDMX) is one of Latin America's top tourist and business travel
destinations. The short-term rental market served by Airbnb spans 16 of the city's
boroughs (\textit{alcaldías}), from the dense historic centre to peripheral
residential areas with very different supply-demand dynamics.

The primary audience for this dashboard is a market analyst at a short-term rental
investment firm or a policy researcher studying housing supply effects. Key questions
the dashboard is designed to answer:

\begin{itemize}
  \item Which boroughs concentrate supply, and which are underserved relative to
        demand signals such as pricing and review scores?
  \item Is CDMX's Airbnb market driven by casual individual hosts or by professional
        property management companies?
  \item What is the price distribution by accommodation type, and where do outliers
        cluster geographically?
\end{itemize}

These questions do not require predictive modeling. They are answered through
descriptive aggregations and interactive charts — the core of a data analyst's
toolkit.

% ---------------------------------------------------------------------------
\section{Data Source}
\label{sec:airbnb-data}
% ---------------------------------------------------------------------------

\sourcefiles{%
  \filepath{projects/00-demo-aestehtics/data-pipeline/airbnb\_etl.py} (lines 23--33)
}

The data comes from \textbf{Inside Airbnb} (\url{http://insideairbnb.com}), an
independent project that scrapes public Airbnb listing pages and publishes city
snapshots under a Creative Commons licence. The CDMX snapshot used here is from
March 2025.

\begin{table}[h]
  \centering
  \caption{Dataset overview — Inside Airbnb CDMX, March 2025}
  \label{tab:airbnb-dataset}
  \begin{tabular}{ll}
    \toprule
    \textbf{Property}     & \textbf{Value} \\
    \midrule
    Raw file              & \filepath{listings.csv.gz} \\
    Total rows            & 27,051 listings \\
    Total columns         & 79 \\
    Columns selected      & 15 (see ETL) \\
    Price null rate       & 12.9\% \\
    Rating null rate      & 12.6\% \\
    Geographic granularity & 16 alcaldías (borough level, not neighbourhood) \\
    Licence               & Creative Commons CC0 \\
    \bottomrule
  \end{tabular}
\end{table}

The null rates in price and rating are structurally expected: new listings have no
reviews, and listings with draft prices or unpublished pricing appear with null
values in the scrape. The pipeline drops nulls for the specific aggregation that
requires the field rather than imputing them — a deliberate choice to avoid
distorting distribution shapes.

\begin{keyinsight}
  Inside Airbnb prices reflect the \emph{published} nightly rate, not the final
  price paid. Cleaning fees, service fees, and taxes can raise the effective
  checkout price by 20–40\%. Any pricing recommendation derived from this data
  should note this limitation explicitly.
\end{keyinsight}

% ---------------------------------------------------------------------------
\section{ETL Pipeline}
\label{sec:airbnb-etl}
% ---------------------------------------------------------------------------

\sourcefiles{%
  \filepath{projects/00-demo-aestehtics/data-pipeline/airbnb\_etl.py}
}

The ETL is a single Python script (\filepath{airbnb\_etl.py}) that reads the
compressed CSV, applies transformations, and writes five JSON files to
\filepath{public/data/airbnb/}. The script uses only \texttt{pandas} and
\texttt{numpy} — no database, no streaming, no orchestration framework. At 27K
rows this is appropriate: pandas loads the file in under a second.

\subsection{Price String Cleaning}

Airbnb exports prices as formatted strings: \texttt{\$1,200.00} or
\texttt{\$14,500.00}. The pipeline strips the currency symbol and thousand
separators before casting to float:

\begin{lstlisting}[style=python, caption={Price cleaning in \filepath{airbnb\_etl.py} (lines 18--20)}]
def clean_price(series: pd.Series) -> pd.Series:
    """Strip $, commas and convert to float."""
    return (series
        .str.replace("$", "", regex=False)
        .str.replace(",", "", regex=False)
        .astype(float))
\end{lstlisting}

A simpler regex like \texttt{[\textbackslash\$,]} would work but is harder to read
at a glance. The two-step replacement makes the intent explicit to any reviewer.

\subsection{Host Segmentation Logic}

Hosts are classified into three segments based on how many listings they operate.
The thresholds are defined in the \texttt{segment()} closure inside
\texttt{build\_host\_segmentation()}:

\begin{lstlisting}[style=python, caption={Host segmentation thresholds (lines 109--113)}]
def segment(n):
    if n == 1:
        return "casual"
    elif n <= 5:
        return "professional"
    return "enterprise"
\end{lstlisting}

\begin{decisionbox}
  \textbf{Why these thresholds?} The 1 / 2--5 / 6+ split maps to observable
  market behaviours: a single-listing host is almost certainly an individual renting
  their own home or spare room. Hosts with 2--5 listings are likely managing
  investment properties part-time. Hosts with 6+ listings operate at a scale that
  requires systematic pricing, availability management, and cleaning coordination
  — this is a commercial operation regardless of legal form.

  The thresholds are not empirically calibrated to CDMX data; they are heuristic
  conventions used by Airbnb researchers. A rigorous segmentation would cluster on
  multiple variables (listing count, price variance, response rate). For a portfolio
  dashboard the heuristic is sufficient and defensible.
\end{decisionbox}

\subsection{Geo Sampling}

The raw dataset has 27K coordinate pairs. Rendering all of them in a browser scatter
chart would produce an illegible overplotted mass and degrade performance. The
pipeline samples 3,000 points at build time:

\begin{lstlisting}[style=python, caption={Geo sampling in \filepath{airbnb\_etl.py} (lines 65--68)}]
def build_geo_heatmap(df: pd.DataFrame, max_points: int = 3000) -> dict:
    valid = df.dropna(subset=["latitude", "longitude", "price"]).copy()
    if len(valid) > max_points:
        valid = valid.sample(max_points, random_state=42)
\end{lstlisting}

The fixed \texttt{random\_state=42} ensures the sample is reproducible between
pipeline runs. The 3,000-point limit is a UX decision: enough density to reveal
the spatial clustering of Cuauhtémoc and Miguel Hidalgo without overwhelming the
browser's SVG renderer.

\subsection{Output Files}

The pipeline produces five JSON files totalling under 500~KB:

\begin{table}[h]
  \centering
  \caption{ETL output files written to \filepath{public/data/airbnb/}}
  \label{tab:airbnb-outputs}
  \begin{tabular}{lll}
    \toprule
    \textbf{File}                     & \textbf{Builder function}        & \textbf{Consumed by} \\
    \midrule
    \filepath{kpis.json}              & \texttt{build\_kpis()}           & KPI bar \\
    \filepath{price\_distribution.json} & \texttt{build\_price\_distribution()} & \texttt{PriceHistogram} \\
    \filepath{geo\_heatmap.json}      & \texttt{build\_geo\_heatmap()}   & \texttt{GeoScatter} \\
    \filepath{neighborhood\_ranking.json} & \texttt{build\_neighborhood\_ranking()} & \texttt{NeighborhoodBar} \\
    \filepath{host\_segmentation.json} & \texttt{build\_host\_segmentation()} & \texttt{HostSegmentation} \\
    \bottomrule
  \end{tabular}
\end{table}

% ---------------------------------------------------------------------------
\section{Architecture}
\label{sec:airbnb-architecture}
% ---------------------------------------------------------------------------

\sourcefiles{%
  \filepath{projects/00-demo-aestehtics/src/app/airbnb/page.tsx},
  \filepath{projects/00-demo-aestehtics/next.config.js},
  \filepath{projects/00-demo-aestehtics/package.json}
}

\subsection{Static JSON Pattern}

The architecture of the Airbnb dashboard is deliberately the simplest possible:
pre-computed JSON files are read at server-render time using Node's \texttt{fs}
module and passed as props to client components. There is no API route, no database
connection, and no runtime data fetching.

\begin{lstlisting}[style=typescript, caption={Static JSON loading in \filepath{src/app/airbnb/page.tsx} (lines 11--14)}]
function readJson<T>(filename: string): T {
  const filePath = join(process.cwd(), 'public', 'data', 'airbnb', filename)
  return JSON.parse(readFileSync(filePath, 'utf-8')) as T
}
\end{lstlisting}

The page component calls \texttt{readJson()} five times — once per dataset — and
forwards all values as typed props to \texttt{AirbnbDashboard}. Because the page
is a React Server Component (no \texttt{'use client'} directive), the file I/O
happens once on the server during the request and the result is serialized into
the initial HTML payload. The client receives fully-hydrated data and never issues
a network request for chart data.

\begin{decisionbox}
  \textbf{Why static JSON instead of an API?} Three reasons:

  \begin{enumerate}
    \item \textbf{Data size}: Five JSON files total under 500~KB. This is small
          enough to embed in a server render without latency penalties.
    \item \textbf{Update frequency}: The Inside Airbnb dataset is a monthly snapshot.
          Data is not changing in real time. An API that queries a database to return
          the same pre-aggregated values on every request adds infrastructure cost
          with zero user-facing benefit.
    \item \textbf{Deployment simplicity}: The dashboard is a static export merged
          into the portfolio's Cloudflare Pages project. No Cloud Run service, no
          database provisioning, no secret management beyond what already exists
          for the rest of the portfolio.
  \end{enumerate}

  The trade-off is that updating the data requires re-running \filepath{airbnb\_etl.py}
  and committing the new JSON files. For a monthly-refresh dataset served to a
  portfolio audience this is an acceptable workflow.
\end{decisionbox}

\subsection{Next.js App Router and Proxy Configuration}

The project uses Next.js~14.2 with the App Router. The route \filepath{/airbnb}
is handled by \filepath{src/app/airbnb/page.tsx}. A separate section of the app
(\filepath{/olist}) connects to a FastAPI backend via a server-side rewrite proxy
defined in \filepath{next.config.js}:

\begin{lstlisting}[style=typescript, caption={Rewrite proxy for the Olist backend in \filepath{next.config.js}}]
async rewrites() {
  const backend = process.env.OLIST_BACKEND_URL || 'http://localhost:2050'
  return [
    {
      source: '/api/olist/:path*',
      destination: `${backend}/:path*`,
    },
  ]
},
\end{lstlisting}

The Airbnb section of the app does not use this proxy at all — it is purely static.
The proxy exists for the Olist (e-commerce) section that shares the same Next.js
process, demonstrating that both patterns (static JSON and live API) can coexist
in the same project.

\subsection{Dependency Choices}

\begin{table}[h]
  \centering
  \caption{Key runtime dependencies and their roles}
  \label{tab:airbnb-deps}
  \begin{tabular}{lll}
    \toprule
    \textbf{Package}      & \textbf{Version} & \textbf{Role} \\
    \midrule
    \texttt{recharts}     & 2.13             & All chart rendering \\
    \texttt{framer-motion} & 11.0            & KPI counter animations \\
    \texttt{lucide-react} & 0.454            & Icon system (theme toggle, back arrow) \\
    \texttt{swr}          & 2.4              & Data fetching for Olist section (not Airbnb) \\
    \texttt{clsx}         & 2.1              & Conditional CSS class composition \\
    \texttt{tailwind-merge} & 2.5            & Conflict-safe Tailwind class merging \\
    \bottomrule
  \end{tabular}
\end{table}

Recharts was chosen over D3.js, Victory, and Nivo. D3 gives maximum control but
requires imperative DOM manipulation that conflicts with React's declarative model
unless carefully managed. Recharts wraps D3's scales in React components, maintaining
the familiar \texttt{<BarChart>} / \texttt{<Scatter>} declarative API while
exposing enough customization hooks (custom shapes, custom tooltips, the
\texttt{Customized} component for SVG overlays) to handle the geo-label overlay
on the scatter chart.

% ---------------------------------------------------------------------------
\section{Visualization Design}
\label{sec:airbnb-viz}
% ---------------------------------------------------------------------------

\sourcefiles{%
  \filepath{projects/00-demo-aestehtics/src/app/globals.css},
  \filepath{projects/00-demo-aestehtics/src/components/charts/},
  \filepath{projects/00-demo-aestehtics/src/components/ui/ThemeToggle.tsx}
}

\subsection{Dark/Light Mode Implementation}

The theme system uses a CSS custom properties (variables) approach rather than
duplicating color values across components. Light and dark tokens are declared in
\filepath{globals.css}:

\begin{lstlisting}[language={},caption={CSS token system in \filepath{src/app/globals.css} (lines 5--47)}]
:root {
  --background: #FAFAF8;   /* warm white, not pure #FFFFFF */
  --foreground: #1A1A1A;
  --chart-tick:   #6B6B6B;
  --chart-label:  #1A1A1A;
  --chart-grid:   #E5E4DF;
  --bar-rank-1:   #1E3A5F; /* dark navy -- top-ranked bars */
  --bar-rank-2:   #2B527F;
  --bar-rank-3:   #9AB0C8; /* lighter blue -- lower-ranked bars */
}

.dark {
  --background: #141414;
  --foreground: #F0EFEB;
  --chart-tick:   #9A9A9A; /* bumped from #6B6B6B for WCAG AA */
  --chart-label:  #F0EFEB;
  --chart-grid:   #2a2a2a;
  --bar-rank-1:   #6BA3C8; /* inverted: lighter in dark mode */
  --bar-rank-2:   #4A7FA8;
  --bar-rank-3:   #9AB0C8;
}
\end{lstlisting}

Chart components reference these tokens via \texttt{var(--chart-tick)} in Recharts
\texttt{tick} and \texttt{label} props. The dark class is toggled on
\texttt{document.documentElement} by \texttt{ThemeToggle.tsx}, which reads the
user's system preference via \texttt{window.matchMedia} on first render and
persists the explicit choice to \texttt{localStorage}.

\subsection{WCAG Contrast Considerations}

The muted text color \texttt{\#6B6B6B} has a contrast ratio of 3.45:1 against the
dark background \texttt{\#141414} — below WCAG AA's 4.5:1 requirement for body
text. The fix is a single CSS override:

\begin{lstlisting}[language={},caption={Muted text contrast fix in \filepath{globals.css} (lines 64--67)}]
/* #6B6B6B gives 3.45:1 on #141414 -- fails AA for small text.
   Override to #9A9A9A which gives 6.6:1. */
.dark .text-muted {
  color: #9A9A9A;
}
\end{lstlisting}

A subtler issue arises with Recharts: SVG \texttt{fill} and \texttt{stroke}
attributes are \emph{not} affected by CSS \texttt{color} inheritance. A hardcoded
\texttt{fill="\#6B6B6B"} in a tick label will fail in dark mode because the SVG
attribute takes precedence over any CSS rule targeting the element. The solution is
to use \texttt{fill='var(--chart-tick)'} in every Recharts \texttt{tick} prop, as
done consistently across all chart components in this project.

\subsection{Chart Type Rationale}

\begin{table}[h]
  \centering
  \caption{Chart type selection rationale for each section}
  \label{tab:airbnb-charts}
  \begin{tabular}{p{3cm}p{3.5cm}p{6cm}}
    \toprule
    \textbf{Chart}         & \textbf{Type}             & \textbf{Rationale} \\
    \midrule
    Price distribution     & Stacked \texttt{BarChart} (histogram) &
      Bins the continuous price variable to reveal the bimodal
      structure (entire home vs.\ private room floors). Stacked bars by room type
      show composition without requiring a separate facet. \\
    \addlinespace
    Geographic distribution & \texttt{ScatterChart} (lon/lat) &
      Recharts \texttt{ScatterChart} maps longitude to x and latitude to y,
      producing an approximate map. A true tile-based map (Leaflet/Mapbox) was
      deferred to avoid token management overhead for a portfolio demo. The
      \texttt{Customized} component overlays borough labels as SVG text. \\
    \addlinespace
    Neighborhood ranking   & Horizontal \texttt{BarChart} &
      Long category names (borough names in Spanish) fit better on a horizontal
      axis with \texttt{layout="vertical"}. Three sort options (listing count,
      avg price, avg rating) are toggled client-side without re-fetching data. \\
    \addlinespace
    Host segmentation      & Vertical \texttt{BarChart} + table &
      Three segments (casual/professional/enterprise) map cleanly to three bars.
      The adjacent summary table surfaces host counts and top operators that a bar
      chart alone cannot convey. \\
    \bottomrule
  \end{tabular}
\end{table}

\subsection{The Geo Scatter Limitation}

Using \texttt{ScatterChart} to approximate a map is a pragmatic compromise with
known failure modes. Recharts maps data coordinates linearly to SVG coordinates,
so the output is an \emph{equirectangular} projection. For CDMX's latitude range
(19.15°N to 19.6°N) this introduces roughly 0.4\% east-west distortion — visually
imperceptible. The main limitation is the absence of base-map tiles: there is no
coastline, no street grid, no borough boundary polygon. Borough labels are drawn
at hardcoded centroid coordinates defined in \filepath{GeoScatter.tsx} (lines
22--38). This is adequate for communicating spatial clustering to a non-specialist
audience but would not survive scrutiny from a GIS professional.

\subsection{Color Encoding on the Geo Scatter}

Point colors encode price tier using a three-level threshold function:

\begin{lstlisting}[style=typescript, caption={\filepath{src/components/charts/GeoScatter.tsx} (lines 47--51)}]
function priceToColor(price: number): string {
  if (price < 800) return '#6B9DB8'   // steel blue -- budget
  if (price < 2000) return '#1E3A5F'  // dark navy -- mid range
  return '#C4841D'                    // amber -- premium
}
\end{lstlisting}

The three-color encoding was chosen deliberately over a continuous gradient. A
continuous sequential scale would require a colorbar legend and demands that viewers
compare gradient shades at small dot sizes — a perceptually difficult task. Three
discrete tiers (budget / mid / premium) allow pre-attentive detection of spatial
patterns: the viewer can immediately see that amber dots cluster along the western
edge without needing to read a colorbar.

The amber accent \texttt{\#C4841D} is the same highlight color used across the
portfolio's chart palette, establishing visual consistency between projects.

% ---------------------------------------------------------------------------
\section{Key Findings}
\label{sec:airbnb-findings}
% ---------------------------------------------------------------------------

\begin{keyinsight}
  \textbf{Cuauhtémoc concentration}: 12,514 listings (46\% of total) sit in a single
  borough. Roma Norte, Condesa, and Centro Histórico within Cuauhtémoc drive this
  cluster. The runner-up borough (Miguel Hidalgo) has approximately half the listing
  count. Tláhuac, at the far southeastern edge of the city, has 40 listings — a
  300$\times$ gap from the leader.
\end{keyinsight}

\begin{keyinsight}
  \textbf{Enterprise host concentration}: The host segmentation reveals a structurally
  professional market. Enterprise hosts (6+ listings) represent only 7\% of all
  hosts yet control 40\% of total supply. The top three enterprise operators are
  Blueground (221 listings), Mr. W (164), and Clau (156). Casual hosts — those with
  a single listing — number 8,370 but collectively hold fewer units than the enterprise
  segment controls per host on average.
\end{keyinsight}

\begin{keyinsight}
  \textbf{Peripheral pricing premium}: Despite far fewer listings, Tlalpan and
  Cuajimalpa show the highest average nightly prices in the dataset (MXN 2,493 and
  MXN 2,151 respectively). This suggests undersupplied premium inventory: listings
  in those boroughs can command high prices because competition is thin. A new host
  entering Tlalpan faces a different competitive environment than one entering
  Cuauhtémoc.
\end{keyinsight}

\begin{keyinsight}
  \textbf{Room type pricing floors}: The price distribution is not unimodal. Entire
  home/apartment listings peak in the MXN 1,000--1,500 bin. Private rooms form a
  secondary mass concentrated near MXN 500. This bimodal shape disappears when the
  two room types are aggregated — a reminder that averaging across heterogeneous
  product categories produces a mean that is representative of neither group.
\end{keyinsight}

Median availability of 16 days per 30 confirms strong demand pressure — listings
are occupied roughly half the time. The 87\% of listings with review scores above
4.5 suggests either genuine quality or a platform selection effect (underperforming
listings are delisted before accumulating reviews).

% ---------------------------------------------------------------------------
\section{Challenge and Interview Questions}
\label{sec:airbnb-questions}
% ---------------------------------------------------------------------------

\begin{challengebox}
  \textbf{Static JSON vs. live API — when does the pattern break?}

  The dashboard reads pre-computed JSON files from the filesystem at request time.
  Identify three scenarios where this architecture would require replacement with a
  live API, and describe what changes to the codebase and infrastructure would be
  needed for each.

  \medskip
  \textit{Consider:} update frequency requirements, personalisation needs, dataset
  size growth, and the cost of adding a database versus the cost of stale data.
\end{challengebox}

\begin{challengebox}
  \textbf{SVG fill inheritance and dark mode}

  A developer adds a new chart component and writes:
  \begin{verbatim}
  <XAxis tick={{ fill: '#6B6B6B' }} />
  \end{verbatim}
  Explain why this breaks in dark mode even though the dark mode CSS class is
  correctly applied to \texttt{document.documentElement}. Describe two ways to fix
  it, and explain which approach is more maintainable in a multi-chart project.
\end{challengebox}

\begin{challengebox}
  \textbf{Geo scatter vs. choropleth map}

  The current geographic visualization plots individual listing coordinates as
  colored dots. An alternative would be a choropleth map where each borough polygon
  is filled with a color proportional to listing count or average price.

  Compare the two approaches on: (a) information density, (b) accessibility for
  colorblind users, (c) implementation complexity given the current Recharts stack,
  and (d) suitability for the three audience questions listed in
  Section~\ref{sec:airbnb-business}.
\end{challengebox}

\begin{interviewbox}
  \textbf{Explain the host segmentation design to a product manager.}

  You have classified hosts into casual (1 listing), professional (2--5 listings),
  and enterprise (6+ listings). A PM asks: ``Why those thresholds? Could the
  segmentation be done better?''

  Give a one-minute answer that acknowledges the heuristic nature of the current
  approach, proposes a data-driven alternative (such as clustering), and explains
  the trade-off between analytical rigor and interpretability for a dashboard
  audience.
\end{interviewbox}

\begin{interviewbox}
  \textbf{Stakeholder question: why are enterprise hosts cheaper than casual hosts?}

  The host segmentation data shows that enterprise hosts have a lower average
  nightly price than casual hosts. A stakeholder finds this counterintuitive and
  asks for an explanation.

  Walk through at least two alternative hypotheses (e.g., volume pricing strategy,
  geographic mix, accommodation type mix) and explain how you would distinguish
  between them using the available data without additional data collection.
\end{interviewbox}
